The following is the equal-case theorem of
\cite[lines~643--656]{DeLasCuevas2017Irreducible}.
The periodic-overlap and coefficient-power hypotheses that earlier statements
carried are now derived from equality of the generated matrix-product vectors,
and the theorem is obtained as a single unconditional result for tensors in
irreducible form II, Theorem~\ref{thm:periodic_ft_equal_case_form_two}. The
source theorem for the unrestricted irreducible form, stated below, still
awaits the passage between the two forms.

% Remaining gap for the source-level statement: the formal equal-case theorem
% assumes trace-preserving block maps (irreducible form II). The passage from
% irreducible form to irreducible form II by a block-diagonal similarity,
% asserted by the source at lines 330--332, is not formalized.
\begin{theorem}[Periodic Fundamental Theorem]\notready
    \label{thm:periodic_ft}
    \uses{def:irreducible_form}
    Let $A$ and $B$ be MPS tensors in irreducible form. Write
    \begin{align}
        A^i & =\bigoplus_{j\in J}(R_j\otimes A_j^i), \notag \\
        B^i & =\bigoplus_{k\in K}(S_k\otimes B_k^i), \notag
    \end{align}
    where $\{A_j\}_{j \in J}$ and $\{B_k\}_{k \in K}$ are bases of periodic
    tensors, $R_j$ and $S_k$ are diagonal multiplicity matrices of sizes
    $r_j\times r_j$ and $s_k\times s_k$, and $A_j^i$ and $B_k^i$ act on
    bond spaces of dimensions $D_j$ and $D_k$.
    Suppose that $\mc{V}(A)=\mc{V}(B)$
    (equal MPV families for all lengths $N$).

    Then $|J|=|K|$ and there is a bijection $\pi:J\to K$
    such that $A_j$ and $B_{\pi(j)}$ are $m_j$-periodic for the same
    period $m_j$.
    Moreover, for each $j\in J$ there exists:
    \begin{enumerate}
        \item an invertible matrix $Y_j$ of size $D_{\pi(j)}\times D_j$,
              and
        \item a diagonal unitary matrix $Z_j$ of size $r_j\times r_j$
              satisfying $Z_j^{m_j}=\Id_{r_j}$,
    \end{enumerate}
    such that $B_{\pi(j)}^i=Y_jA_j^iY_j^{-1}$ for every physical
    index~$i$, and, after reordering the diagonal entries of the
    multiplicity matrix $S_{\pi(j)}$, one has $Z_jR_j=S_{\pi(j)}$.
    Equivalently, setting $Z=\bigoplus_j Z_j\otimes\Id_{D_j}$ gives an
    invertible matrix $Y$ such that $ZA^i=YB^iY^{-1}$ for all $i$,
    with $[A^i,Z]=0$ and $\mc{V}(A)=\mc{V}(ZA)$.

    \noindent\emph{Remark.}
    Theorem~\ref{thm:periodic_ft_equal_case_form_two} is this statement for
    tensors in irreducible form II, that is, with trace-preserving block maps.
    What separates the two is the passage between irreducible form and
    irreducible form II by a block-diagonal similarity, asserted by the source
    at lines 330--332 and recorded in
    \cite{gap:dccsp17_periodic_overlap_route_alignment}, section ``Scope
    restriction: the equal case in irreducible form II''.
\end{theorem}

\begin{theorem}[Diagonal multiplicity gauge from equal powers]
    \label{thm:zgauge_construction}
    \lean{MPSTensor.zgauge_construction}
    \leanok
    \uses{lem:z_gauge_diagonal_pow, lem:z_gauge_diagonal_mul}
    Let $(\mu_i)_{i\in I}$ and $(\nu_i)_{i\in I}$ be two finite lists of
    nonzero multiplicity entries. If $\mu_i^m=\nu_i^m$ for every $i\in I$,
    then there exists a diagonal matrix $Z$ such that
    \begin{align}
        Z^m           & =\Id, \notag        \\
        Z\,\diag(\nu) & =\diag(\mu). \notag
    \end{align}
\end{theorem}

\begin{proof}\leanok
    Take $Z=\diag(\mu_i/\nu_i)_{i\in I}$ and apply the two diagonal
    identities from Lemmas~\ref{lem:z_gauge_diagonal_pow} and
    \ref{lem:z_gauge_diagonal_mul}.
\end{proof}

\begin{theorem}[Scalar multiplicity gauge from power sums]
    \label{thm:equalcase_zgauge_power_sums}
    \lean{MPSTensor.equalCase_zgauge_of_power_sums}
    \leanok
    \uses{thm:zgauge_construction}
    Let $(\mu_i)_{i\in I}$ and $(\nu_i)_{i\in I}$ be two finite lists of
    multiplicity entries. Suppose that the $\nu_i$ are nonzero, that
    $\mu_i^m=\nu_i^m$ for every $i\in I$, and that
    $\sum_{i\in I}\mu_i^k=\sum_{i\in I}\nu_i^k$ for every positive
    integer~$k$. Then the two lists determine the same multiset, and there is
    a diagonal $Z$ satisfying $Z^m=\Id$ and
    $Z\,\diag(\nu)=\diag(\mu)$.
\end{theorem}

\begin{proof}\leanok
    Apply Theorem~\ref{thm:zgauge_construction} to obtain $Z$, and apply the
    Newton--Girard power-sum identity to recover the multiplicity multiset.
\end{proof}

\begin{theorem}[Conditional matching of bases of periodic tensors]
    \label{thm:periodic_equalcase_matching}
    \lean{MPSTensor.fundamentalTheorem_periodic_equalCase_matching}
    \leanok
    \uses{def:irreducible_form}
    Let $A$ and $B$ be in irreducible form, and suppose their bases of periodic
    tensors contain no repeated pair of distinct basis elements. If the
    periodic-overlap hypothesis holds for the two bases, then the two bases
    have the same cardinality and are matched by a bijection whose matched
    elements are repeated blocks.
\end{theorem}

\begin{proof}\leanok
    Apply the conditional block-matching theorem to the two bases extracted
    from the irreducible-form representations.
\end{proof}

\begin{theorem}[Conditional scalar equal-case component]
    \label{thm:periodic_equalcase_scalar_component}
    \lean{MPSTensor.fundamentalTheorem_periodic_equalCase}
    \leanok
    \uses{thm:periodic_equalcase_matching, thm:equalcase_zgauge_power_sums}
    Let $A$ and $B$ be in irreducible form, and suppose their bases of periodic
    tensors contain no repeated pair of distinct basis elements. Assume the
    periodic-overlap hypothesis and the corresponding power equality for the
    scalar multiplicity entries. Then the bases are matched by a bijection, and
    each matched pair has a scalar diagonal $Z_j\in\MN{1}$ with
    \begin{align}
        Z_j^{m_j}             & =\Id, \notag                   \\
        Z_j\,\diag(\mu_{A,j}) & =\diag(\mu_{B,\pi(j)}). \notag
    \end{align}
    The two scalar multiplicity entries also determine the same singleton
    multiset.
\end{theorem}

\begin{proof}\leanok
    First use Theorem~\ref{thm:periodic_equalcase_matching} to match the
    bases of periodic tensors. For each matched block, apply
    Theorem~\ref{thm:equalcase_zgauge_power_sums} to the singleton list of
    multiplicity entries.
\end{proof}

\begin{lemma}[Power weights of a sector decomposition]%
    \label{lem:sector_power_weights}
    \lean{MPSTensor.SectorDecomposition.powWeights,
        MPSTensor.SectorDecomposition.coeff_powWeights}
    \leanok
    \uses{def:sector_weight_data}
    Raising every multiplicity entry of a sector decomposition to a fixed
    power $m$ leaves the basis blocks and the multiplicities unchanged, and
    the coefficient of the resulting decomposition at length $n$ is the
    original coefficient at length $mn$:
    \begin{align}
        \sum_{q}\bigl(\mu_{j,q}^{m}\bigr)^{n}
        =\sum_{q}\mu_{j,q}^{mn}.
        \notag
    \end{align}
\end{lemma}

\begin{proof}\leanok
    The entries $\mu_{j,q}^{m}$ are again nonzero, so they form a sector
    weight family over the same basis, and the displayed identity is the
    exponent rule.
\end{proof}

\begin{lemma}[Matched multiplicity entries along a period progression]%
    \label{lem:matched_sector_weight_pow_equiv_period}
    \lean{MPSTensor.matched_sector_weight_pow_equiv_of_period_multiple}
    \leanok
    \uses{lem:sector_power_weights, lem:sector_bnt_matched_sector_weight_equiv}
    Let two sector decompositions have distinguished blocks $j_0$ and $k_0$,
    with multiplicity entries $(\mu_{j_0,q})_q$ and $(\nu_{k_0,q})_q$, and let
    $\zeta$ be a nonzero scalar. Suppose that
    \begin{align}
        \sum_{q}\mu_{j_0,q}^{mn}
        =\zeta^{mn}\sum_{q}\nu_{k_0,q}^{mn}
        \notag
    \end{align}
    for every $n$ beyond some threshold. Then the two blocks carry the same
    multiplicity, and there is a bijection $\tau$ between their copy indices
    with
    \begin{align}
        \nu_{k_0,\tau(q)}^{m}=\bigl(\zeta^{-1}\mu_{j_0,q}\bigr)^{m}
        \notag
    \end{align}
    for every $q$. This is the form in which the power-sum comparison enters
    \cite[lines~681--684]{DeLasCuevas2017Irreducible}, where the comparison is
    available only along the multiples of a period; the recovery of the
    entries from their power sums is
    \cite[lines~537--569]{DeLasCuevas2017Irreducible}.
\end{lemma}

\begin{proof}\leanok
    Raise all multiplicity entries of both decompositions to the power $m$.
    By Lemma~\ref{lem:sector_power_weights} the hypothesis becomes a
    comparison of the coefficients of the two resulting decompositions at
    every exponent beyond the threshold, with scalar $\zeta^{m}$. Applying
    Lemma~\ref{lem:sector_bnt_matched_sector_weight_equiv} to those
    decompositions gives the multiplicity equality and the bijection, and its
    conclusion is the displayed identity.
\end{proof}

\begin{theorem}[Equal-case coefficient extraction]%
    \label{thm:periodic_equalcase_coefficient_extraction}
    \lean{MPSTensor.SectorDecomposition.coeff_eq_of_sameMPV_of_matched_basis}
    \leanok
    \uses{def:same_mpv2_pos, def:periodic_tensor, lem:paper_decomp_formula,
        thm:periodic_basis_nonzero_subfamily_li,
        thm:pgvwc07_periodic_state_vector_decomposition}
    Let
    \begin{align}
        A^i & =\bigoplus_{j\in J}(R_j\otimes A_j^i),
        \notag                                       \\
        B^i & =\bigoplus_{k\in K}(S_k\otimes B_k^i),
        \notag
    \end{align}
    be two multiplicity-bearing decompositions whose bases of periodic
    tensors are matched by a bijection $\pi\colon J\to K$ and nonzero scalars
    $\zeta_j$ with
    \begin{align}
        \ket{V^{(N)}(A_j)}=\zeta_j^{N}\ket{V^{(N)}(B_{\pi(j)})}
        \notag
    \end{align}
    for every length $N$, with $\{A_j\}_{j\in J}$ pairwise non-repeated and
    $A_j$ periodic of period $m_j$. If $\mc{V}^+(A)=\mc{V}^+(B)$, then there is
    an $N_0$ such that
    \begin{align}
        \tr\bigl(R_j^{N}\bigr)
        =\zeta_j^{-N}\,\tr\bigl(S_{\pi(j)}^{N}\bigr)
        \notag
    \end{align}
    for every $N\ge N_0$ with $m_j\mid N$ and every $j\in J$.
    This is \cite[lines~672--680]{DeLasCuevas2017Irreducible}. The scalar
    $\zeta_j$ is displayed rather than suppressed: the identity compares the
    power sums of $R_j$ with the \emph{rescaled} power sums of $S_{\pi(j)}$,
    which is what the source means at
    \cite[lines~667--671]{DeLasCuevas2017Irreducible} by absorbing the phase
    into $S_j$.
\end{theorem}

\begin{proof}\leanok
    Expand both vectors along the basis $\{A_j\}_{j\in J}$ by
    Lemma~\ref{lem:paper_decomp_formula}, rewriting each
    $\ket{V^{(N)}(B_{\pi(j)})}$ as $\zeta_j^{-N}\ket{V^{(N)}(A_j)}$. Equality
    of the two vectors makes the combination of the $\ket{V^{(N)}(A_j)}$ with
    coefficients $\tr(R_j^N)-\zeta_j^{-N}\tr(S_{\pi(j)}^N)$ vanish. By
    Theorem~\ref{thm:pgvwc07_periodic_state_vector_decomposition} the term
    for an index with $m_j\nmid N$ contributes the zero vector, so the
    combination runs over the indices with $m_j\mid N$, and by
    Theorem~\ref{thm:periodic_basis_nonzero_subfamily_li} those vectors are
    independent for every large $N$. Each coefficient therefore vanishes.
\end{proof}

\begin{theorem}[Blockwise multiplicity gauge in the equal case]%
    \label{thm:periodic_equalcase_blockwise_zgauge}
    \lean{MPSTensor.equalCase_blockwise_zgauge}
    \leanok
    \uses{thm:periodic_equalcase_coefficient_extraction,
        lem:matched_sector_weight_pow_equiv_period, thm:zgauge_construction}
    Under the hypotheses of
    Theorem~\ref{thm:periodic_equalcase_coefficient_extraction}, for every
    $j\in J$ the multiplicity matrices $R_j$ and $S_{\pi(j)}$ have the same
    size $r_j$, and there are a reordering $\tau$ of the diagonal entries of
    $S_{\pi(j)}$ and a diagonal matrix $Z_j$ with
    \begin{align}
        Z_j^{m_j}=\Id_{r_j},\qquad
        Z_j\,\bigl(\zeta_j R_j\bigr)=T_jS_{\pi(j)}T_j^{\dagger},
        \notag
    \end{align}
    where $T_j$ implements $\tau$. This is
    \cite[lines~681--688]{DeLasCuevas2017Irreducible}. The scalar $\zeta_j$
    stays visible: the diagonal entries of $Z_j$ are $m_j$-th roots of unity
    against the rescaled multiplicity matrix $\zeta_jR_j$, and the rescaling
    is in general not trivial even for period one, as the pair of blocks
    $B_j=e^{i\theta}A_j$ shows.
\end{theorem}

\begin{proof}\leanok
    Fix $j$. Theorem~\ref{thm:periodic_equalcase_coefficient_extraction}
    compares the power sums of $R_j$ and $S_{\pi(j)}$ at every multiple of
    $m_j$ beyond $N_0$, with scalar $\zeta_j^{-1}$. Feeding that progression
    to Lemma~\ref{lem:matched_sector_weight_pow_equiv_period} gives the
    equality of sizes, the reordering $\tau$, and the entry identity
    $\nu_{\pi(j),\tau(q)}^{m_j}=(\zeta_j\mu_{j,q})^{m_j}$. Since the entries
    $\zeta_j\mu_{j,q}$ are nonzero,
    Theorem~\ref{thm:zgauge_construction} produces the diagonal $Z_j$ with the
    two displayed properties.
\end{proof}

\begin{theorem}[Global multiplicity gauge and similarity]%
    \label{thm:periodic_equalcase_global_zgauge}
    \lean{MPSTensor.equalCase_global_zgauge_of_blockwise}
    \leanok
    \uses{def:periodic_tensor, def:same_mpv2_pos,
        thm:pgvwc07_periodic_state_vector_decomposition}
    Let
    \begin{align}
        A^i & =\bigoplus_{j\in J}(R_j\otimes A_j^i),
        \notag                                       \\
        B^i & =\bigoplus_{k\in K}(S_k\otimes B_k^i),
        \notag
    \end{align}
    be two multiplicity-bearing decompositions, let $\pi\colon J\to K$ be a
    bijection with $D_j=D_{\pi(j)}$, let $\tau_j$ reorder the multiplicity
    copies of the matched blocks, let $\xi_j$ be scalars and $Y_j$ invertible
    matrices with
    \begin{align}
        A_j^i=\xi_jY_jB_{\pi(j)}^iY_j^{-1},
        \notag
    \end{align}
    and let $z_{j,l}$ be scalars with
    \begin{align}
        z_{j,l}\,\xi_j\mu_{j,l}=\nu_{\pi(j),\tau_j(l)},
        \qquad z_{j,l}^{m_j}=1,
        \notag
    \end{align}
    where $A_j$ is $m_j$-periodic and every $m_j$ divides $L$. Writing $Z_j$
    for the diagonal matrix of the $z_{j,l}$ and
    $Z=\bigoplus_{j\in J}Z_j\otimes\Id_{D_j}$, one has $Z^L=\Id$,
    $[A^i,Z]=0$ for every $i$, $\mc{V}^+(A)=\mc{V}^+(ZA)$, and there is an
    invertible $Y$ with $ZA^i=YB^iY^{-1}$ for every $i$.

    This is the assembly of
    \cite[lines~689--690]{DeLasCuevas2017Irreducible}, with $Y$ the direct sum
    $\bigoplus_jT_j\otimes Y_j$ of the copy reorderings and the blockwise
    similarities.
\end{theorem}

\begin{proof}\leanok
    The multiplicity gauge acts by the scalar $z_{j,l}$ on the bond space of
    the copy $(j,l)$, so it commutes with every block-diagonal matrix and its
    $L$-th power is the identity. Multiplying $A^i$ by it rescales the
    multiplicity entries by the $z_{j,l}$, and the hypothesis on the
    $z_{j,l}$ makes the rescaled entries equal to the matched entries of $B$.
    The matched blocks are therefore related copy by copy, so the two
    assembled tensors are related by the permutation of the matched copies
    followed by the direct sum of the blockwise similarities. For the
    matrix-product vectors, a copy whose period divides the length contributes
    the factor $z_{j,l}^N=1$, and a copy whose period does not divide the
    length contributes the zero vector by
    Theorem~\ref{thm:pgvwc07_periodic_state_vector_decomposition}.
\end{proof}

\begin{theorem}[Fundamental theorem for MPS in irreducible form II, equal case]%
    \label{thm:periodic_ft_equal_case_form_two}
    \lean{MPSTensor.fundamentalTheorem_periodic_equalCase_sectorDecomposition}
    \leanok
    \uses{def:periodic_tensor, def:same_mpv2_pos,
        thm:periodic_equalcase_global_zgauge,
        thm:periodic_equalcase_blockwise_zgauge,
        thm:periodic_equalcase_coefficient_extraction,
        thm:periodic_conditional_block_matching,
        thm:periodic_multiplicity_nondecaying_overlap,
        thm:zgauge_construction}
    Let $A$ and $B$ carry the multiplicity-bearing forms
    \begin{align}
        A^i & =\bigoplus_{j\in J}(R_j\otimes A_j^i),
        \notag                                       \\
        B^i & =\bigoplus_{k\in K}(S_k\otimes B_k^i),
        \notag
    \end{align}
    where $R_j$ and $S_k$ are diagonal with nonzero entries and the two
    displayed families are pairwise non-repeated bases of periodic tensors
    whose blocks are trace preserving, of periods $m_j$ and $n_k$. If
    $\mc{V}^+(A)=\mc{V}^+(B)$, then there is a bijection
    $\pi\colon J\to K$ such that, for every $j$, one has $m_j=n_{\pi(j)}$,
    the blocks $A_j$ and $B_{\pi(j)}$ are repeated with unit-modulus scalar
    $\xi_j$, the multiplicity matrices $R_j$ and $S_{\pi(j)}$ have the same
    size, and there are a reordering $T_j$ of the diagonal entries of
    $S_{\pi(j)}$ and a diagonal matrix $Z_j$ with
    \begin{align}
        Z_j^{m_j}=\Id,\qquad
        Z_j\left(\xi_jR_j\right)=T_jS_{\pi(j)}T_j^\dagger.
        \notag
    \end{align}
    Writing $Z=\bigoplus_{j\in J}Z_j\otimes\Id_{D_j}$, one has
    $Z^{\operatorname{lcm}_j m_j}=\Id$, $[A^i,Z]=0$ for every $i$,
    $\mc{V}^+(A)=\mc{V}^+(ZA)$, and there is an invertible $Y$ with
    $ZA^i=YB^iY^{-1}$ for every $i$.

    \emph{Scope restriction.}  This is
    \cite[Theorem~3.8, lines~643--693]{DeLasCuevas2017Irreducible} for tensors
    in irreducible form II, that is, with trace-preserving block maps. The
    source theorem assumes only irreducible form and asserts at lines 330--332
    that one passes between the two forms by a block-diagonal similarity, so
    that a proof in either form applies to the other; that passage is not
    carried out here. The restriction and its elimination plan are recorded in
    \cite{gap:dccsp17_periodic_overlap_route_alignment}, section ``Scope
    restriction: the equal case in irreducible form II''. The unrestricted
    source theorem is Theorem~\ref{thm:periodic_ft}, which remains
    unformalized.

    The unit-modulus scalar of each matched pair is displayed rather than
    suppressed. The entries of $Z_j$ are $m_j$-th roots of unity against the
    rescaled multiplicity matrix $\xi_jR_j$, which is what the source means at
    lines 667--671 by absorbing the phase into $S_j$; the scalar is genuinely
    present, as the period-one pair $B_j=e^{i\theta}A_j$ shows.
\end{theorem}

\begin{proof}\leanok
    Equality of the generated matrix-product vectors is in particular
    proportionality, so
    Theorem~\ref{thm:periodic_multiplicity_nondecaying_overlap} supplies
    non-decaying partners and
    Theorem~\ref{thm:periodic_conditional_block_matching} matches the two
    bases; matched blocks have equal periods because repeated blocks have the
    same peripheral spectrum. Each matched pair carries a unit-modulus scalar
    $\xi_j$ relating its matrix-product vectors at every length, so
    Theorem~\ref{thm:periodic_equalcase_coefficient_extraction} compares the
    multiplicity power sums at every large multiple of $m_j$, and
    Theorem~\ref{thm:periodic_equalcase_blockwise_zgauge} turns that
    comparison into the reordering $T_j$ and the entry identities
    $\nu_{\pi(j),\tau_j(l)}^{m_j}=(\xi_j\mu_{j,l})^{m_j}$. The ratios
    $\nu_{\pi(j),\tau_j(l)}/(\xi_j\mu_{j,l})$ are therefore $m_j$-th roots of
    unity, and Theorem~\ref{thm:zgauge_construction} collects them into
    $Z_j$. Theorem~\ref{thm:periodic_equalcase_global_zgauge} assembles the
    blockwise data into the global $Z$ and $Y$.
\end{proof}
